\chapter{REVUE DE LITTERATURE}
\thispagestyle{empty}

\section{Recommandation fondee sur les graphes dans Wikipedia}

Wikipedia se modele naturellement comme un graphe dans lequel les noeuds sont
des articles et les aretes sont des hyperliens. Cette structure en fait un bon
terrain d'experimentation pour la recommandation fondee sur les graphes et
l'apprentissage de representations. Wiki-CS~\cite{mernyei2020wikics} fournit
l'un des sous-ensembles de reference les plus connus pour les experiences de
reseaux de neurones sur graphes a partir de contenu derive de Wikipedia, mais
il s'agit avant tout d'un benchmark de classification de noeuds plutot que
d'un systeme de recommandation a proprement parler.

La recommandation sur Wikipedia a egalement ete etudiee plus directement.
WikiRecNet a propose un cadre de recommandation a grande echelle utilisant la
structure de Wikipedia pour la recuperation d'articles~\cite{wikirecnet2020}.
Une lecon importante de cette ligne de travaux est que la topologie du graphe
contient un signal fort de proximite, mais que le probleme de recommandation ne
se confond pas avec une prediction generique sur graphe. Ce point est
particulierement important ici, puisque la tache cible est la recommandation
d'article a article et non la classification de noeuds.

\section{Plongements textuels semantiques}

Les methodes de plongement textuel fournissent une seconde famille de signaux.
Les modeles de langue fondes sur les Transformers~\cite{vaswani2017attention},
tels que BERT~\cite{devlin2019bert}, ont rendu pratique l'encodage de phrases et de
documents sous forme de vecteurs denses preservant les relations semantiques.
Dans la litterature de recherche par similarite, des travaux tels que
Sentence-BERT~\cite{reimers2019sentencebert} ont egalement joue un role
important en montrant comment des architectures de type BERT peuvent etre
adaptees a la comparaison efficace de segments textuels dans un espace dense.
Des modeles plus recents, orientes vers la recuperation, ameliorent encore
cette capacite pour la recherche par similarite. BGE-M3, en particulier, est
concu comme un retriever dense multilingue et couvre plusieurs usages
orientes recuperation~\cite{chen2025bgem3}. Cela en fait un candidat solide
pour la recommandation d'articles a partir du texte brut de Wikipedia.

Dans un cadre de recommandation, les embeddings semantiques ont deux avantages.
Premierement, ils peuvent retrouver des sujets lies meme lorsque la structure
de liens est peu dense. Deuxiemement, ils produisent souvent des voisinages
facilement interpretables par un humain, car les elements recuperes sont relies
par le sens textuel et non seulement par leur position dans le graphe. Leur
limite principale est qu'ils peuvent devenir tres dependants de la popularite
ou trop larges semantiquement lorsqu'ils sont utilises seuls, surtout sur de
grandes collections d'articles.

\section{Plongements structurels}

Les embeddings structurels cherchent a capter la position dans le graphe, le
contexte de voisinage local et la structure communautaire directement a partir
des aretes. Node2Vec est l'un des exemples les plus influents de cette
famille~\cite{grover2016node2vec}. En combinant des marches aleatoires avec un
apprentissage de type skip-gram, il apprend des representations vectorielles
dans lesquelles des noeuds ayant des contextes structurels semblables deviennent
proches dans l'espace des embeddings.

Pour la recommandation, les embeddings structurels sont attractifs parce qu'ils
capturent des signaux que le texte peut manquer. Deux articles peuvent etre peu
semblables dans leur formulation tout en occupant des positions presque
identiques dans le graphe. Les methodes structurelles obtiennent ainsi souvent
de tres bons resultats sur la recuperation de liens caches. Leur limite est
qu'elles peuvent sur-accentuer les hubs, les effets de degre et certains
artefacts propres au graphe si le protocole d'evaluation n'est pas
suffisamment controle.

\section{Reseaux de neurones sur graphes}

Les reseaux de neurones sur graphes generalisent l'apprentissage de
representations structurelles en propageant l'information a travers les aretes.
GCN~\cite{kipf2017semi} est l'une des architectures les plus utilisees et reste
une baseline transductive solide. GraphSAGE~\cite{hamilton2017inductive}
etend cette idee par agregation de voisinage et est souvent utilise lorsque le
comportement inductif importe. GAT introduit un mecanisme d'attention sur les
voisinages du graphe~\cite{velickovic2018graph}, ce qui en fait un modele
prometteur lorsque l'importance des aretes est susceptible de varier fortement
selon les voisins.

La litterature sur la recommandation comprend egalement des systemes hybrides
qui combinent modeles textuels et modeles de graphe. Les etudes de synthese sur
la recommandation multimodale~\cite{zhou2023multimodal} et des exemples
recents de systemes hybrides fondes sur BERT et GNN~\cite{anime_hybrid_gnn_bert}
montrent que les canaux d'information se comportent rarement de maniere
identique et que la conception de la fusion importe beaucoup. De petits choix
de ponderation, de normalisation ou d'alignement des ensembles candidats
peuvent modifier significativement le classement final. Ce point est
particulierement important dans ce projet, ou la concatenation naive, la
concatenation corrigee, les hybrides avec pooling et la fusion par rang ne
conduisent pas aux memes conclusions.

\section{Les LLM comme couches de presentation}

Les grands modeles de langue sont de plus en plus utilises dans les pipelines
de recommandation, mais leur role peut varier fortement. Dans certains
systemes, ils generent des explications; dans d'autres, ils raffinent le
classement; ailleurs encore, ils reorganisent des elements deja recuperes dans
des structures plus lisibles pour l'humain. Une partie recente de la
litterature etudie les LLM comme recommandeurs ou rerankers a part entiere,
par exemple en classement zero-shot~\cite{hou2023zeroshotrankers}, en
instruction tuning pour le top-k~\cite{luo2023recranker} ou dans des cadres
d'analyse plus larges sur l'usage des LLM en recommandation~\cite{xu2024tapping}.
Le projet actuel n'emploie toutefois les LLM que dans un role plus etroit. Le
LLM ne recupere pas de nouveaux articles, ne produit pas de contenu factuel et
ne definit pas le benchmark. Il reorganise au contraire une liste fixe de
titres deja recuperes dans une presentation de type parcours d'apprentissage.

Cette conception est importante pour la fiabilite. Elle contraint le LLM a
travailler sur un ensemble ferme de candidats et rend plus simple la validation
ou le rejet de sa sortie. Dans le cadre du projet, cette couche a evolue
d'experiences initiales avec de petits modeles locaux vers un schema de sortie
structuree plus controle, accompagne d'une validation cote application. Ce
positionnement se rapproche davantage de la litterature recente sur la
generation structuree et la conformite a des schemas que d'un usage du LLM
comme moteur principal de recommandation~\cite{geng2025structured}.

\section{Positionnement du present travail}

Ce projet se situe a l'intersection de l'exploration de graphes, de la
recuperation semantique et de l'ingenierie de prototypes de recommandation. Par
rapport aux travaux existants, ses traits distinctifs sont les suivants:
\begin{itemize}
    \item l'utilisation d'un graphe Wikipedia d'informatique enrichi avec le
    texte reel des articles plutot qu'avec de simples identifiants de noeuds;
    \item la comparaison de baselines semantiques, structurelles, hybrides et
    fondees sur des GNN dans un meme cadre experimental;
    \item un accent explicite mis sur la difference entre la recuperation de
    liens et la qualite de recommandation au sommet de la liste classee;
    \item la separation entre la couche de recuperation et la couche optionnelle
    d'organisation par LLM.
\end{itemize}

La contribution finale mise en oeuvre n'est donc pas un modele unique et
entierement nouveau, mais un pipeline experimental soigneusement documente qui
montre comment la conception methodologique, la strategie de fusion et le point
de vue adopte pour l'evaluation peuvent changer l'interpretation des resultats
de recommandation sur Wikipedia.
